\documentclass[10pt]{beamer}

\usepackage[utf8]{inputenc}
\usepackage[spanish]{babel}
\usepackage{graphicx}

\mode<presentation>
\usetheme{Madrid}
%\usecolortheme[RGB={111,73,135}]{structure}
\usecolortheme[RGB={128,0,0}]{structure}
%\usecolortheme[RGB={0,96,0}]{structure}
%\usecolortheme[RGB={200,0,200}]{structure}
%\usecolortheme[RGB={0,128,0}]{structure}
%\usecolortheme[RGB={0,0,128}]{structure}
\usefonttheme{serif}
\useinnertheme{rectangles}
\useoutertheme{split}
\setbeamercovered{transparent}


% Definiciones para usar luego
\newtheorem{ejemplo}{Ejemplo}
\newtheorem{definicion}{Definición}

\title{Entendiendo el Kernel de Linux}
\author{Jesús Espino García\\Samuel Rodríguez Sevilla}
\date{13 de Marzo de 2009}
\subject{Entendiendo el Kernel de Linux}

\institute[GUL UC3M]{
  Grupo de Usuarios de Linux\\
  Universidad Carlos III de Madrid.\\
  \ \\
  \includegraphics[height=1.5cm]{imgs/gul}}

\setcounter{tocdepth}{2}

\AtBeginSubsection[]
{
  \begin{frame}<beamer>{Indice}
    \tableofcontents[sectionstyle=show/hide,subsectionstyle=show/shaded/hide]
  \end{frame}
}

\begin{document}

  \frame{\maketitle}
  \begin{frame}
	\setcounter{tocdepth}{1}
  	\tableofcontents[sectionstyle=show/show,subsectionstyle=show/shaded/hide]
	\setcounter{tocdepth}{2}
  \end{frame}
  
  \section{Introducción}
\subsection{¿Que es?}
\begin{frame}
	\begin{itemize}
		\item Núcleo del sistema operativo GNU/Linux.
		\item GPL.
		\item Estable.
		\item Maduro.
		\item Moderno.
		\item Eficiente.
		\item Versátil.
		\item Configurable.
		\item y mucho mas!!
	\end{itemize}
\end{frame}

\subsection{Un poco de historia}
\begin{frame}
	\begin{itemize}
		\item En Abril de 1991 Linus Torvals empieza el desarrollo de Linux.
		\item En Septiembre del 1991 lanza la versión 0.01
		\item En Diciembre del 1991 libera la primera versión ''self hosted''.
		\item En Marzo del 1994 aparece la versión 1.0.0.
		\item En 1996 aparece Tux como mascota de Linux y la versión 2.0 del mismo.
		\item En 1999 lanza la versión 2.2
		\item En 2001 lanza la versión 2.4
		\item En 2003 lanza la versión 2.6 (la actual).
	\end{itemize}
\end{frame}

\subsection{Numeración de versiones}
\begin{frame}
	\begin{block}{Hasta la versión 2.5}
		\begin{itemize}
			\item El primer numero indicaba la versión de la rama del kernel.
			\item El segundo indicaba si era estable o no (par estable, impar inestable).
			\item El tercero era la minor-release de esa rama.
		\end{itemize}
	\end{block}
	\begin{block}{versión 2.6.0 y posteriores}
		\begin{itemize}
			\item El primer y segundo numero se han estabilizado (no se esperan cambios a corto plazo).
			\item El tercer numero indica la versión del kernel.
			\item Si esta versión esta seguida por -rcX significa que es una versión de desarrollo.
			\item Si esta versión esta seguida por un .X significa que es estable y que esta en el bugfix X.
		\end{itemize}
	\end{block}
\end{frame}

\subsection{Git}
\begin{frame}{Git}
	\begin{itemize}
		\item Sistema de control de versiones del kernel de Linux.
		\item Distribuido.
		\item Desarrollado directamente por Linus Torvals.
		\item Apareció como una alternativa a BitKeeper.
		\item Existen muchas ramas desarrolladas de manera distribuida.
	\end{itemize}
\end{frame}

\subsection{Visión General}
\begin{frame}{A vista de pájaro}
	\begin{center}
		\includegraphics[height=6.5cm]{imgs/vista1}
	\end{center}
\end{frame}

\begin{frame}{Capas}
	\begin{center}
		\includegraphics[height=6.5cm]{imgs/vista2}
	\end{center}
\end{frame}

\begin{frame}{Subsistemas}
	\begin{center}
		\includegraphics[height=6.5cm]{imgs/vista3}
	\end{center}
\end{frame}

  \section{Subsistema de Memoria}
\subsection{Direcciones de memoria}
\begin{frame}{Tipos de Direcciones}
	\begin{itemize}
		\item Dirección lógica: Dirección del segmento mas el offset.
		\item Dirección lineal (o virtual): Dirección en toda la memoria.
		\item Dirección física: Dirección de acceso a la RAM.
	\end{itemize}
\end{frame}

\subsubsection{Segmentos}
\begin{frame}{¿Qué es un segmento?}
	\begin{itemize}
		\item Conjunto de memoria.
		\item Se utiliza para establecer unos permisos a la memoria.
		\item Se accede a sus contenidos a través de dirección del segmento + offset.
		\item Esta implementado en hardware.
	\end{itemize}
\end{frame}

\begin{frame}{¿Cómo lo usa Linux?}
	\begin{itemize}
		\item Lo usa de un modo muy limitado.
		\item Se prefiere el uso de páginas.
		\item Linux crea 4 segmentos que abarca toda la memoria (cada uno).
		\begin{itemize}
			\item usercode
			\item userdata
			\item kernelcode
			\item kerneldata
		\end{itemize}
		\item Como todos los segmentos empiezan en 0, la posición se define únicamente por el offset.
		\item Linux solo usa segmentos cuando lo requiere la arquitectura.
	\end{itemize}
\end{frame}

\subsubsection{Paginas}
\begin{frame}{¿Qué son las páginas?}
	\begin{itemize}
		\item Conjunto de memoria.
		\item Se utiliza para establecer unos permisos a la memoria.
		\item Se accede a su contenido a través de una dirección compuesta de varios elementos.
		\item Los elementos de la dirección designan posiciones en unas tablas de índices.
		\item Esta implementado en hardware.
	\end{itemize}
\end{frame}

\begin{frame}{¿Cómo las usa Linux?}
	\begin{itemize}
		\item Se usa un modelo con 4 tablas de direccionamiento (para poder albergar mucha memoria):
		\begin{itemize}
			\item Page Global Directory
			\item Page Upper Directory
			\item Page Middle Directory
			\item Page Table
		\end{itemize}
		\item Se usa el mismo modelo para 32 y 64 bits.
		\item Cuando se usa 32 bits y no se usa PAE, Linux elimina las tablas Upper y Middle.
		\item Cuando se usa 32 bits y se usa PAE, Linux elimina la tabla Upper.
	\end{itemize}
\end{frame}

\subsection{Gestión de memoria}
\subsubsection{Buddy System}
\begin{frame}{Buddy System}
	\begin{itemize}
		\item Sistema encargado de reservar las páginas y hacer un seguimiento de las páginas libres.
		\item Soluciona el problema de la fragmentación externa.
		\item Se almacenan las páginas en 11 listas dependiendo del grupo de páginas contiguas disponibles.
		\item Linux utiliza varios Buddy System, uno para cada ''zona'' de memoria (DMA, normal, high).
	\end{itemize}
\end{frame}
\subsubsection{Slab Allocator}
\begin{frame}{Slab Allocator}
	\begin{itemize}
		\item El Buddy System reserva memoria como mínimo del tamaño de la página.
		\item El Slab Allocator permite reservar bloques de tamaños arbitrarios mas pequeños que una página.
		\item Funciona encima del Buddy System.
		\item Tiene un sistema de caches que permite reutilizar bloques de memoria.
		\item La reutilización mejora el rendimiento de Linux que utiliza los slab con frecuencia.
	\end{itemize}
\end{frame}

\subsubsection{Reclamado de páginas}
\begin{frame}{Reclamado de páginas}
	\begin{itemize}
		\item Al hacer una petición al Buddy System y no encontrar páginas apropiadas, este hace un reclamado.
		\item El sistema de reclamado divide las páginas en 4 tipos:
		\begin{itemize}
			\item Unreclaimable: Paginas que no pueden o no es necesario reclamar.
			\item Swappable: Paginas que pueden pasarse a la swap.
			\item Syncable: Paginas que pueden ser sincronizadas a disco.
			\item Discardables: Paginas que pueden ser descargadas.
		\end{itemize}
		\item Dependiendo del tipo de páginas hace una acción u otra.
		\item El reclamado de páginas también se ejecuta periódicamente mediante kswapd.
	\end{itemize}
\end{frame}

  \section{Interrupciones}
\subsection{Tipos interrupciones}
\begin{frame}{Tipos de interrupciones}
	\begin{itemize}
		\item Excepciones: Interrupciones del propio procesador.
		\item Interrupciones: Señales (normalmente por IRQ) del hardware.
	\end{itemize}
\end{frame}

\subsection{IRQs y APIC}
\begin{frame}{IRQs y APIC}
	\begin{block}{IRQ}
		\begin{itemize}
			\item IRQ es una señal especial que va directamente a la CPU.
			\item Tiene su propio canal de comunicación (normalmente incluido en el bus del sistema).
			\item Este bus permite mandar señales al procesador en cualquier momento.
		\end{itemize}
	\end{block}
	\begin{block}{APCI}
		\begin{itemize}
			\item APIC es un sistema de IRQs en 2 niveles (normalmente) para sistemas SMP.
			\item Permite discriminar el envío de IRQs a un procesador u otro.
			\item Cada CPU tiene su sistema de APIC local para gestionar sus interrupciones.
		\end{itemize}
	\end{block}
\end{frame}

\subsection{Excepciones}
\begin{frame}{Excepciones}
	\begin{itemize}
		\item Indica un error emitido por la CPU.
		\item Normalmente es una situación anómala producida por la programación.
		\item Ocurre durante la ejecución del proceso y no hay cambio de contexto.
	\end{itemize}
\end{frame}

\subsection{Interrupciones}
\begin{frame}{Interrupciones}
	\begin{itemize}
		\item Indica una solicitud al procesador de que se haga algo.
		\item Normalmente las interrupciones las producen el hardware.
		\item Puede ser urgente o no.
		\item En caso de ser urgente se interrumpe la ejecución y lanza el manejador.
		\item En caso de no ser urgente se espera al momento idóneo para lanzar el manejador.
		\item Para manejar una interrupción se produce un cambio de contexto.
		\item Se pueden emitir interrupciones por software (softirq).
		\item Las interrupciones pueden ser interrumpidas por otras.
	\end{itemize}
\end{frame}

  \section{Subsistema de procesos}
\subsection{Los procesos}
\subsubsection{¿Qué es un proceso?}
\begin{frame}{¿Qué es un proceso?}
	\begin{itemize}
		\item Estructura de datos en el SO.
		\item Proceso completo o hilo.
		\item Instancia de tarea a realizar por la CPU.
		\item Utiliza, reserva y comparte recursos.
	\end{itemize}
\end{frame}

\subsubsection{Tipos de procesos}
\begin{frame}{Tipos de procesos}
	\begin{itemize}
		\item Procesos: Instancia de un programa en ejecución.
		\item Procesos ligeros: Instancia del ''contador de programa''.
	\end{itemize}
\end{frame}

\subsubsection{Descriptor de proceso}
\begin{frame}{Descriptor de proceso}
	\begin{block}{Información que contiene}
		\begin{itemize}
			\item Estado.
			\item Memoria asignadas.
			\item Ficheros abiertos.
			\item Directorio actual.
			\item Padre.
			\item Señales recibidas.
			\item y mucho más.
		\end{itemize}
	\end{block}
\end{frame}

\subsubsection{Estados}
\begin{frame}{Estados de un proceso}
	\begin{itemize}
		\item TASK\_RUNNING: Proceso en la CPU o a espera de CPU.
		\item TASK\_INTERRUPTIBLE: Proceso durmiendo a espera de algo.
		\item TASK\_UNINTERRUPTIBLE: Como TASK\_INTERRUPTIBLE, pero sin poder ser interrumpido.
		\item TASK\_STOPPED: Se ha parado la ejecución del proceso.
		\item TASK\_TRACED: Se ha parado la ejecución del proceso por un debugger.
		\item TASK\_ZOMBIE: Se ha terminado el proceso pero no hay ningún wait esperándole.
		\item TASK\_DEAD: El proceso esta siendo eliminado definitivamente.
	\end{itemize}
\end{frame}

\subsubsection{Cambio de proceso}
\begin{frame}{Cambio de proceso}
	\begin{itemize}
		\item Muchos cambios de proceso por segundo.
		\item Actividad necesariamente rápida.
		\item Se deben a las interrupciones.
		\item La interrupción por tiempo permite hacer planificación de procesos.
		\item Consiste en almacenar el estado del proceso actual y sustituirlo por otro.
		\item El kernel (casi en su totalidad) es ''reentrant''.
	\end{itemize}
\end{frame}

\subsubsection{Crear procesos}
\begin{frame}{Crear procesos}
	\begin{block}{Procesos}
		\begin{itemize}
			\item Proceso 0: El swapper o idle, se crea al inicio de manera ''artesanal''.
			\item Proceso 1: El init, se crea como copia del proceso 0 y ejecuta el programa init (ya en modo usuario).
			\item Resto de los procesos: Se crean como copia del proceso padre (normalmente init).
		\end{itemize}
	\end{block}
	\begin{block}{Estrategias}
		\begin{itemize}
			\item Copy On Write: Padre e hijo comparten memoria para lectura y en el momento de la escritura se reserva una nueva pagina.
			\item Procesos ligeros: Padre e hijo comparte memoria.
		\end{itemize}
	\end{block}
\end{frame}

\subsubsection{Destruir procesos}
\begin{frame}{Destruir procesos}
	\begin{itemize}
		\item Recibe la syscall \_exit().
		\item Pasa del estado en que este a TASK\_ZOMBIE.
		\item Su padre hace un wait, y pasa al estado TASK\_DEAD.
		\item Se eliminan todos los datos del proceso.
		\item Se elimina su descriptor del sistema.
		\item Los compiladores normalmente incluyen \_exit() de manera automática al final del binario.
	\end{itemize}
\end{frame}

\subsubsection{Hilos del kernel}
\begin{frame}{Hilos del kernel}
	\begin{itemize}
		\item Son procesos creados dentro del kernel.
		\item Algunos se crean al inicio.
		\item Otros bajo demanda.
		\item Algunos ejemplos son:
		\begin{itemize}
			\item kapmd: Gestiona eventos relacionados con APM.
			\item kswapd: Reclama memoria periódicamente.
			\item pdflush: Vacía los buffers ''sucios'' a disco.
		\end{itemize}
	\end{itemize}
\end{frame}

\subsection{Planificación (Scheduling)}
\begin{frame}{Algoritmos}
	\begin{block}{O(1)}
		\begin{itemize}
			\item Scheduler hasta el 2.6.22
			\item Basado en listas de prioridad.
			\item Calculo del quamto basado en prioridad estática.
			\item Calculo del sucesor basado en prioridad dinámica.
			\item La prioridad estática viene dada por el tipo de proceso (normal o tiempo real) y el nice.
			\item La prioridad dinámica se calcula a partir de la prioridad estática y el ''average sleep'' del proceso.
		\end{itemize}
	\end{block}
\end{frame}

\begin{frame}{Algoritmos}
	\begin{block}{CFS (Completely Fair Scheduler)}
		\begin{itemize}
			\item Scheduler desde 2.6.23
			\item Mas simple
			\item Reparte el tiempo del procesador de manera equitativa.
			\item Puntúa mas a los procesos que han usado menos el procesador.
			\item Penaliza a los que más.
			\item Intenta que el reparto sea justo.
		\end{itemize}
	\end{block}
\end{frame}

\subsection{Memoria de procesos}
\begin{frame}{Memoria de procesos}
	\begin{itemize}
		\item Un proceso usar un conjunto de direcciones (regiones de memoria).
		\item Cada proceso cuenta con una pila que puede usar para albergar memoria dinámica.
		\item Cada región de memoria tiene unos permisos de acceso para el proceso.
		\item Varios procesos pueden acceder a las mismas posiciones de memoria (bibliotecas, shared memory\dots).
		\item El intento de acceso a cualquier otra posición de memoria genera una excepción.
		\item El intento de uso no debido de las posiciones de memoria asignadas genera una excepción.
		\item Al acceder a su memoria puede ser que no este la pagina (por diversos motivos), por lo cual se produce una demanda de pagina.
	\end{itemize}
\end{frame}

  \section{syscalls y señales}
\subsection{syscalls}
\begin{frame}{syscalls}
	\begin{itemize}
		\item Son el mecanismo que tienen los procesos para comunicarse con el núcleo.
		\item Todas devuelven un valor de tipo \emph{long} para indicar si ha habido éxito.
		\item Están identificadas por un valor que utiliza la aplicación cuando le quiere indicar al núcleo qué \emph{syscall} desea utilizar.
		\item Todas se encuentran en una tabla denominada \emph{sys\_call\_table} que es dependiente de la arquitectura. La tabla suele estar en el fichero \emph{entry.S} y en los i386 se encuentra en \emph{arch/i386/kernel}.
	\end{itemize}
\end{frame}

\begin{frame}{syscalls}
	\begin{itemize}
		\item Se utiliza una interrupción software para indicarle al núcleo que se va a hacer uso de una \emph{syscall}. En la arquitectura i386 se utiliza la interrupción 0x80.
		\item Las \emph{syscalls} tienen parámetros, que en el caso de la arquitectura i386 se pasan por registro. Si hay más de 5 parámetros, éstos serán pasados a través de un buffer de memoria. Núnca se debe acceder a dicho buffer directamente, se copia con \emph{copy\_from\_user} a espacio de núcleo y si se le tiene que pasar información se usa \emph{copy\_to\_user}.	
	\end{itemize}
\end{frame}

\begin{frame}[containsverbatim]{syscalls}
	El formato de una \emph{syscall} es el siguiente:
	\begin{verbatim}
		asmlinkage long sys_nombre(/*parametros*/)
		{
		   return 0;
		}
	\end{verbatim}
	Además se registran en \emph{asm/unistd.h} con un define de formato \_\_NR\_nombre. Y con \emph{\_syscall}$n$ se crea una función que hará la llamada a la \emph{syscall}.
\end{frame}

\subsection{Señales}
\begin{frame}{Señales}
	\begin{itemize}
		\item Las señales son mensajes que se envían a los procesos o grupos de procesos para informarles de ciertos eventos, como alarmas, fallos de segmentación de memoria, etc.
		\item Es un sistema bastante complejo internamente que requiere de más de 6 tipos diferentes de estructuras.
		\item Se puede obtener una lista de las señales en \emph{man 7 signal}.
		\item Hay dos señales que no pueden ser capturadas por los procesos: SIGKILL y SIGSTOP.
		\item El núcleo divide la gestión de señales en dos partes:
		\begin{itemize}
			\item Generación de señales, dónde se actualizan las tablas del proceso para representar que se ha enviado la señal, y
			\item Despacho de señales, dónde se fuerza al proceso a reaccionar ante la señal.
		\end{itemize}
		\item El núcleo comprueba constantemente el estado de las señales.
	\end{itemize}
\end{frame}
\begin{frame}{Señales}
	Los procesos pueden realizar alguna de las siguientes acciones con las señales:
	\begin{itemize}
		\item Ignorarlas
		\item Ejecutar una de las acciones por defecto, que son:
		\begin{itemize}
			\item Terminar, que hace que el proceso no siga ejecutándose.
			\item Realizar un volcado de memoria (generar un fichero \emph{core}) y luego termina al proceso.
			\item Ignorar.
			\item Parar, que deja al procesos suspendido hasta que reciba una señal.
			\item Continuar, que en caso de que el proceso se encontrase parado lo reanuda.
		\end{itemize}
		\item Capturar la señal y tratarla con un manejador de señal.
	\end{itemize}

\end{frame}

%\begin{frame}{Generación de señales}
%\end{frame}
%\begin{frame}{Descpaho de señales}
%\end{frame}

  \section{Input/Output}
\subsection{Dispositivos}
\begin{frame}{Dispositivos}
	\begin{itemize}
		\item Los dispositivos son todos los elementos de un ordenador que permiten realizar operaciones de entrada/salida. Por ejemplo los discos duros, el teclado, etc.
		\item Los dispositivos se representan en \emph{sysfs} con un sistema de \emph{kobjects} (se encuentra declarado en \emph{linux/kobject.h}). Estos \emph{kobjects} describen la información sobre un dispositivo para que puede ser fácilmente consultable.
		\item Para poder acceder a los dispositivos se necesitan los ficheros de dispositivo (típicamente se encuentran en \emph{/dev/}). Estos vienen descritos por dos números denominados \emph{major} y \emph{minor}  y por su tipo, que puede ser caracter o bloque.
	\end{itemize}
\end{frame}

\begin{frame}{Dispositivos}
	\begin{itemize}
		\item El \emph{major} identifica el tipo de dispositivo y el \emph{minor} identifica un dispositivo concreto de ese tipo.
		\item Los dispositivos permiten realizar al programador una serie de acciones (\emph{open}, \emph{read}, \emph{lseek}, \emph{ioctl}, etc.).
		\item Los controladores de dispositivo deben estar registrados en el sistema y utilizan un descriptor de tipo \emph{device\_driver} que debe encontrarse en dentro de una estructura de tipo \emph{device}.
		\item Los controladores de dispositivo compilados estáticamente en el núcleo se registran durante el proceso de inicio. El resto deberán llamar a la función de registro que le corresponda (por ejemplo \emph{pci\_register\_driver}).
		\item Casi todos los controladores de dispositivo siguen una política de asignación de recursos tardía. De este modo, asignan la memoria cuando se les solicita hacer algo (normalmente por una llamada a \emph{open}) y la liberan cuando ya no es necesaria (no hay ninguna aplicación utilizando sus recursos).
	\end{itemize}
\end{frame}

\subsection{Arquitectura de I/O en bloques}
\subsubsection{Arquitectura}
\begin{frame}{Arquitectura}
	Hay dos tipos de dispositivos:
	\begin{itemize}
		\item Dispositivos de caracter: leen y/o escriben un flujo de datos. Su acceso es secuencial. Ejemplos de este tipo de dispositivos son el ratón y el teclado.
		\item Dispositivos de bloque: leen y/o escriben bloques de datos y permiten el acceso aleatorio a los mismos. Ejemplos de este tipo de dispositivos tenemos los discos duros, DVDs, etc.
	\end{itemize}
\end{frame}

\begin{frame}{Arquitectura}
	Los dispositivos de bloques entrañan más complejidad.
	\begin{itemize}
		\item Los dispositivos se dividen en sectores.
		\item Los sectores se agrupan en bloques.
		\item Los bloques se agrupan en segmentos (un segmento es la cantidad de información que puede mandarse por un canal DMA).
		\item Los segmentos se agrupan en páginas (pueden tener bloques que no se encuentran en un segmento).
	\end{itemize}
\end{frame}

\begin{frame}{Arquitectura}
	\begin{itemize}
		\item Los bloques se leen a memoria en áreas denominadas buffers.
		\item Los buffers está definidos por la estructura \emph{buffer\_head} que se encuentra en \emph{linux/buffer\_head.h}.
		\item Esta estructura controla los datos sobre a qué bloque corresponde el buffer, de qué dispositivo, en qué estado se encuentra, cuantas referencias hay a él, etc.
	\end{itemize}
\end{frame}

\begin{frame}{Arquitectura}
	\begin{itemize}
		\item Las operaciones de I/O que se están realizando se encuentran definidas por la estructura \emph{bio} que se encuentra en \emph{linux/bio.h}.
		\item Esta estructura representa las operaciones sobre páginas.
		\item La operación de I/O puede afectar a más de una página. Se controla con el array \emph{bio\_vec} de \emph{bio}.
		\item Las operaciones se encolan un una cola de peticiones pendientes.
	\end{itemize}
\end{frame}

\subsubsection{IO Scheduler}
\begin{frame}{IO Scheduler}
	\begin{itemize}
		\item Las operaciones de I/O no son óptimas en tiempo.
		\item El sistema operativo planifica estas operaciones para mejorar el rendimiento.
		\item La planificación se lleva a cabo siguiendo alguna de las siguientes políticas (en el núcleo se denominan \emph{elevators}):
		\begin{itemize}
			\item \emph{Linus Elevator}
			\item \emph{Deadline I/O Scheduler}
			\item \emph{Anticipatory I/O Scheduler}
			\item \emph{Complete Fair Queuing I/O Scheduler}
			\item \emph{Noop I/O Scheduler}
		\end{itemize}
	\end{itemize}
\end{frame}

\subsection{VFS}

  \include{linux_vfs}
  \section{Inicio del sistema}
\subsection{Inicio del sistema}
\begin{frame}{Inicio del sistema}
	\begin{itemize}
		\item Es el proceso que permite poner en funcionamiento al sistema operativo.
		\item Comienza por la BIOS, que realiza una comprobación del hardware y pasa el control a un programa que típicamente se encuentra en el sector de inicio del disco duro.
		\item El cargador de arranque se encarga de cargar el núcleo en memoria y pasarle el control junto con una serie de parámetros (se llama a la función \emph{setup( )} que se encuentra desplazada 0x200 bytes del inicio del núcleo).
	\end{itemize}
\end{frame}

\subsection{Proceso de inicio}
\begin{frame}{Proceso de inicio}
	El proceso de arranque realiza las siguientes tareas:
	\begin{itemize}
		\item Se invoca a la BIOS para saber la cantidad de RAM del sistema o, en sistemas ACPI, obtener un mapa físico de la RAM.
		\item Ajusta la la velocidad del repetición del teclado.
		\item Inicializa el adaptador de vídeo.
		\item Reinicia el controlador de disco y obtiene los parámetros de los discos.
		\item Comprueba que haya un bus IBM Micro Channel.
		\item Comprueba si hay un ratón PS/2.
		\item Comprueba si la BIOS soporta APM (\emph{Advanced Power Management}).
		\item En caso de que la BIOS soporte EDD (\emph{Enhanced Disk Drive Services}), se solicita crear una tabla en memoria con la descripción de los discos duros.
	\end{itemize}
\end{frame}
\begin{frame}{Proceso de inicio}
	\begin{itemize}
 		\item Según si el núcleo era sin comprimir o comprimido, procede a continuar o desplaza el núcleo en memoria para colocarlo donde le corresponde.
		\item Se activa el pin A20 en el controlador de teclado 8042.
		\item Se crea una tabla de descriptores de interrupciones (IDT) provisional y una tabla global de descripción (GDT) provisional.
		\item Reinicia la unidad de coma flotante (FPU).
		\item Se reprograma el controlador programable de interrupciones (PIC) para enmascarar todas las interrupciones menos la de IRQ2.
		\item Se pone la CPU en modo protegido.
		\item Se llama a la función \emph{startup\_32( )}.
	\end{itemize}
\end{frame}

\subsection{startup\_32}
\begin{frame}{La función startup\_32 (I)}
	\begin{itemize}
		\item Inicializa los registros de segmentación y una pila provisional.
		\item Pone a 0 los bits del registro \emph{eflags}.
		\item Inicializa con ceros el área de datos del núcleo que se encuentran en \emph{\_edata} y \emph{\_end}.
		\item Llama a \emph{decompress\_kernel( )}.
		\item Una vez descomprimido el núcleo se llama a un nuevo método \emph{startup\_32~(~)}.
	\end{itemize}
\end{frame}
\begin{frame}{La función startup\_32 (y II)}
	El nuevo \emph{startup\_32( )} hace las siguiente operaciones:
	\begin{itemize}
		\item Inicializa los registros de segmentación con los valores finales.
		\item Rellena el segmento \emph{bss} del núcleo con ceros.
		\item Inicializa la tabla de páginas provisional del núcleo.
		\item Almacena la dirección del \emph{Page Global Directory} en el registro \emph{cr3}.
		\item Asigna la pila de modo núcleo al proceso 0.
		\item Vuelve a poner a 0 el registro \emph{eflags}.
		\item Llama a \emph{setup\_idt( )}.
		\item Almacena los parámetros obtenidos de la BIOS en el primer \emph{frame} de memoria.
		\item Identifica el modelo del procesador.
		\item Carga los registros de la tabla GDT e IDT.
		\item salta a la función \emph{start\_kernel( )}. 
	\end{itemize}
\end{frame}

\subsection{start\_kernel}
\begin{frame}{La función start\_kernel}
	Llama a las siguientes funciones:
	\begin{itemize}
		\item \emph{sched\_init( )}.
		\item \emph{build\_all\_zonelists( )}.
		\item \emph{page\_alloc\_init( )} y \emph{mem\_init( )}.
		\item \emph{trap\_init( )}.
		\item \emph{softirq\_init( )}.
		\item \emph{time\_init( )}.
		\item \emph{kmem\_cache\_init( )}.
		\item \emph{calibrate\_delay( )}.
		\item \emph{kernel\_thread( )}.
	\end{itemize}
\end{frame}


  \section{Para terminar.}

\begin{frame}[containsverbatim]{Referencias}
  \begin{itemize}
    \item \verb+/usr/src/linux/Documentation+: Documentación del kernel.
    \item http://www.kernelnewbies.org: Pagina para iniciarse en el mundo del kernel de Linux.
    \item http://www.kernel.org: Pagina oficial del kernel de Linux.
    \item http://www.makelinux.net/kernel\_map: Mapa interactivo del kernel.
  \end{itemize}
\end{frame}

\begin{frame}{Dudas}
  \dots
\end{frame}

\begin{frame}
  \frametitle{\begin{center}Fin\end{center}}
  \begin{center}
    \includegraphics[height=3cm]{imgs/gul}
  \end{center}
\end{frame}


\end{document}
